% =========================
% sec_kinetics_B.tex — ПОЛНАЯ ПЕРЕЗАПИСЬ ДЛЯ ВИТРИНЫ (B/Математика)
%
% Витринная логика:
% 1) не подменять глобальную тороидальную геометрию локальным патчем;
% 2) показать, зачем вводится Σ_loc;
% 3) сохранить канон: дрейф/вращение только через анизотропию Δp на Σ;
% 4) показать эволюцию языка: от q-схемы к дискретно-модовой схеме m_A, m_S, S1/S2/S4;
% 5) вернуть локальные формулы к глобальному интегралу по Σ.
%
% Требует в преамбуле:
% \usepackage{tikz}
% \usepackage{booktabs}
% \usepackage{array}
% \usetikzlibrary{arrows.meta,calc,positioning,decorations.pathmorphing}
%
% Канон-опора:
% - CANON.md — единственный источник истины;
% - PDF/LaTeX — витрина и может отставать;
% - тонкий тор r << R допускает только локальное квазицилиндрическое приближение;
% - шнек — динамический режим ГП;
% - КУБК = баланс числа контактов, а не источник дрейфа;
% - дрейф/вращение только через анизотропию Δp на Σ;
% - прокси (K,G,A_ij) — ярлыки измерения;
% - e_∥ — ось §17, не h_A и не h_S.
% =========================

\section{Кинетика ЭВО: дрейф и вращение как интегральный отклик анизотропии контактов на \(\Sigma\)}

\subsection{Глобальная цель и временное сужение задачи}

Цель текущего шага --- не заменить глобальную тороидальную геометрию ЭВО локальным цилиндром, а временно сузить задачу до локального патча \(\Sigma_{\mathrm{loc}}\), чтобы явнее выделить кинетику контактов на ГП, ответственную за дрейф и ротацию. Глобально ЭВО остаётся тороидальной или двуториальной структурой; локальный патч используется только как язык записи и вычисления для той же самой канонической геометрии.

Витринно важно сохранить последовательность рассуждения. Сначала задаётся глобальная геометрия и полный интеграл по \(\Sigma\). Затем объясняется, почему в режиме тонкого тора возникает естественная локальная карта \(\Sigma_{\mathrm{loc}}\). После этого становится видно, какие механизмы удобнее отделяются на патче, почему одной непрерывной \(q\)-геометризации оказалось недостаточно, почему пришлось различить \(h_A\) и \(h_S\), и почему шнек теперь следует читать через дискретно-согласованную схему \((m_A,m_S)\) с базовыми модами \(S1/S2/S4\). Финальный шаг --- возврат к глобальному интегралу по \(\Sigma\), чтобы локальный язык не подменял онтологию объекта.

\subsection{Тонкий тор как базовая геометрия и причины локального патча}

\subsubsection{Граничная поверхность \(\Sigma\)}

Пусть \(\Sigma\subset\mathbb{R}^3\) --- компактная кусочно-\(C^2\) ориентируемая поверхность без края
(или конечное объединение таких компонент для многоканального ЭВО). Задана параметризация
\[
X:U\subset\mathbb{R}^2\to\Sigma,\qquad x=X(u,v).
\]
Элемент площади и единичная нормаль:
\[
dS=\|X_u\times X_v\|\,dudv,\qquad
n(x)=\frac{X_u\times X_v}{\|X_u\times X_v\|}.
\]

\subsubsection{Ось \(e_\parallel\) (норма §17)}

Ось \(e_\parallel\) --- \textbf{ось симметрии объекта} и \textbf{ось тестов §17}:
\[
V_\parallel := V\cdot e_\parallel.
\]
Везде далее \(e_\parallel\) используется как нормативная ``χ-ручка''. Важно: \(e_\parallel\) \emph{не} подменяется ни директором траектории Амеров \(h_A\), ни директором шнековой моды \(h_S\).

\subsubsection{Хиральность \(\chi\)}

Фиксируем паспортную хиральность:
\[
\chi\in\{-1,+1\}.
\]

\subsubsection{Контактная статистика на \(\Sigma\) и КУБК}

На \(\Sigma\) необходимо жёстко различать:
\begin{itemize}
\item баланс \emph{числа контактов} --- КУБК;
\item баланс \emph{приращений импульса} --- источник дрейфа и момента.
\end{itemize}

Пусть \(f_{\mathrm{хаос}}(x,v,t)\) и \(f_{\mathrm{канал}}(x,v,t)\) --- локальные распределения скоростей амеров
в окрестности \(\Sigma\) со стороны хаотической фазы и со стороны канала. Нормаль \(n(x)\) направлена из ЭВО наружу.

Интенсивности контактов:
\[
\Phi^-_{\mathrm{хаос}}(x,t)=\int_{v\cdot n\le 0}|v\cdot n|\;f_{\mathrm{хаос}}(x,v,t)\,d^3v,\qquad
\Phi^+_{\mathrm{канал}}(x,t)=\int_{v\cdot n\ge 0}(v\cdot n)\;f_{\mathrm{канал}}(x,v,t)\,d^3v.
\]

\(\Phi^\pm\) трактуются только как \textbf{интенсивности касаний \(\Sigma\)}, а не как ``потоки СЭ''.

КУБК фиксирует отсутствие устойчивого перекоса по \emph{числу} контактов:
\[
\overline{\Phi^-_{\mathrm{хаос}}}(x)=\overline{\Phi^+_{\mathrm{канал}}}(x),\qquad x\in\Sigma.
\]

Это означает: устойчивый дрейф нельзя выводить из перекоса числа контактов. Он может возникать только через анизотропию \(\Delta p\) на \(\Sigma\).

\subsubsection{Полный интеграл по \(\Sigma\)}

Локальная импульсная анизотропия задаётся полем
\[
j(x,t):=\lambda(x,t)\,\mathbb{E}[\Delta p\mid x,t],\qquad x\in\Sigma.
\]
Разложение:
\[
j=j_n\,n+j_\top,\qquad j_n=(j\cdot n),\qquad j_\top\perp n.
\]

Интегральные отклики:
\[
V(t)=\mu_T\int_{\Sigma} j(x,t)\,dS,
\]
\[
\Omega(t)=\tilde\mu_R\int_{\Sigma}\bigl(r(x)-r_0\bigr)\times j(x,t)\,dS.
\]

Именно этот глобальный интеграл по \(\Sigma\) остаётся главным объектом кинетики. Все локальные записи ниже должны возвращаться к нему.

\subsection{Квазицилиндрический патч как локальный язык, а не новая онтология}

\subsubsection{Почему нужен локальный патч}

В режиме тонкого тора \(r\ll R\) изменение локальной кривизны по большой окружности мало на масштабе длины свободного пробега \(l\). Поэтому короткий участок \(\Sigma\) можно записывать как квазицилиндрический патч:
\[
\Sigma_{\mathrm{loc}}\subset\Sigma.
\]

Это не новая геометрия объекта и не альтернатива тору. Это локальная карта, в которой проще отделить:
\begin{itemize}
\item контактную статистику на ГП;
\item касательные компоненты \(j(x)\);
\item локальные tor/sp-вклады;
\item различие между траекторией Амеров и шнековой модой ГП.
\end{itemize}

\subsubsection{Локальный базис патча}

На \(\Sigma_{\mathrm{loc}}\) вводим:
\[
(e_z,e_\vartheta,n),
\]
где
\begin{itemize}
\item \(e_z\) --- единичный касательный вектор вдоль большой окружности тора;
\item \(e_\vartheta\) --- единичный касательный вектор по малой окружности;
\item \(n\) --- внешняя единичная нормаль.
\end{itemize}

Локальное соответствие глобальной тороидальной геометрии:
\[
e_z \leftrightarrow t_\varphi,\qquad e_\vartheta \leftrightarrow t_\theta.
\]

Таким образом, патч --- это лишь локальный язык записи той же граничной поверхности \(\Sigma\).

\subsubsection{Прокси как язык наблюдения}

Сохраняем канонический прокси-стандарт:
\[
K=n\langle v^2\rangle,\qquad
G=n\langle v\rangle,\qquad
P_{ij}=n\langle v_i v_j\rangle,\qquad
A_{ij}=P_{ij}-\frac13\mathrm{tr}(P)\delta_{ij}.
\]
Проекции на касательную плоскость:
\[
\Pi_\top(\xi)=\xi-(\xi\cdot n)n,\qquad
g_\top:=\Pi_\top(G),\qquad
a_\top:=\Pi_\top(A\,n).
\]

Прокси не являются сущностями и не несут причинности; они служат только языком измерения контактной статистики и анизотропии \(\Delta p\).

% =========================================================
% ВСТАВКА 1: БОЛЬШАЯ ВИЗУАЛЬНАЯ СХЕМА
% =========================================================

\subsubsection{Схема локального квазицилиндрического патча и направлений кинетики}

На рис.~\ref{fig:quasi_cyl_patch_scheme} показана расчётная схема локального квазицилиндрического приближения. Глобально ЭВО сохраняет тороидальную геометрию с граничной поверхностью \(\Sigma\). В режиме тонкого тора \(r\ll R\) выделяется локальный патч \(\Sigma_{\mathrm{loc}}\), на котором вводятся координаты \((z,\vartheta)\), внешняя нормаль \(n\), касательные направления \(e_z\) и \(e_\vartheta\), а также два различных касательных директора: \(h_A\) для траекторий Амеров в канале и \(h_S\) для шнековой моды ГП. Важно, что поток Амеров и бег шнека не отождествляются: первый связан с \(h_A\), второй --- с \(h_S\). Локально вычисляются контактные интенсивности, прокси и внутренние вклады \(j_{\mathrm{tor}}^{\mathrm{loc}}\), \(j_{\mathrm{sp}}^{\mathrm{loc}}\), после чего результат собирается обратно в глобальный интеграл по \(\Sigma\).

% =========================================================
% ПОЛНОСТРАНИЧНАЯ ВСТАВКА:
% крупная схема квазицилиндрического патча
%
% Требует в преамбуле:
% \usepackage{tikz}
% \usepackage{graphicx}
% \usetikzlibrary{arrows.meta,calc,positioning,decorations.pathmorphing}
% =========================================================

% =========================================================
% ПОЛНОСТРАНИЧНАЯ ВЕРТИКАЛЬНАЯ СХЕМА:
% рисунки расположены один над другим
%
% Требует в преамбуле:
% \usepackage{tikz}
% \usepackage{graphicx}
% \usetikzlibrary{arrows.meta,calc,positioning,decorations.pathmorphing}
% =========================================================

% =========================================================
% ПОЛНОСТРАНИЧНАЯ ВЕРТИКАЛЬНАЯ СХЕМА
% Готовый блок с уменьшенными шрифтами
% =========================================================

% =========================================================
% ПОЛНОСТРАНИЧНАЯ ВЕРТИКАЛЬНАЯ СХЕМА
% Готовый блок с уменьшенными шрифтами
% =========================================================

% =========================================================
% ОТДЕЛЬНЫЙ ЛИСТ С РИСУНКОМ НА ВСЮ ПЛОЩАДЬ
% Без caption, чтобы схема заняла максимум места
% =========================================================

% =========================================================
% ПОЛНОСТРАНИЧНАЯ СХЕМА — СИЛЬНО ОБЛЕГЧЁННАЯ
% Внутри рисунка только короткие метки
% =========================================================

% =========================================================
% НОВАЯ СХЕМА С НУЛЯ
% Полностраничный рисунок, минимальный текст внутри
%
% Требует в преамбуле:
% \usepackage{tikz}
% \usepackage{graphicx}
% \usepackage{geometry}
% \usetikzlibrary{arrows.meta,calc,positioning}
% =========================================================

% =========================================================
% РИСУНОК А — ГЕОМЕТРИЯ КВАЗИЦИЛИНДРИЧЕСКОГО ПАТЧА
% Полностраничный, без перегруза формулами
%
% Требует в преамбуле:
% \usepackage{tikz}
% \usepackage{geometry}
% \usetikzlibrary{arrows.meta,calc,positioning}
% =========================================================

\clearpage
\begingroup
\thispagestyle{empty}
\newgeometry{top=5mm,bottom=5mm,left=5mm,right=5mm}

\phantomsection
\refstepcounter{figure}
\label{fig:quasi_cyl_geometry}
\addcontentsline{lof}{figure}{Геометрия локального квазицилиндрического патча}

\begin{center}
\resizebox{0.99\textwidth}{0.96\textheight}{%
\begin{tikzpicture}[
    >=Latex,
    font=\Large,
    boundary/.style={line width=1.3pt},
    axis/.style={->, line width=1.25pt},
    amer/.style={->, line width=1.3pt},
    screw/.style={->, dashed, line width=1.3pt},
    lbl/.style={font=\Large},
    sml/.style={font=\large},
    every node/.style={inner sep=2pt}
]

% ---------------------------------------------------------
% Верх: глобальный тор
% ---------------------------------------------------------

\node[lbl] at (-3.8,11.1) {\(\Sigma\)};

\begin{scope}[shift={(0,8.25)}]
    \draw[boundary] (0,0) ellipse (3.35 and 2.0);
    \draw[boundary] (0,0) ellipse (1.58 and 0.95);

    \draw[line width=1.8pt] (-1.15,0.72)
        arc[start angle=145,end angle=35,x radius=1.58,y radius=0.95];

    \draw[rounded corners=5pt, line width=1.2pt] (0.95,0.38) rectangle (2.45,1.28);
    \node[sml] at (3.55,1.15) {\(\Sigma_{\mathrm{loc}}\)};

    \draw[axis] (0,0) -- (2.8,0) node[midway, below] {\(R\)};
    \draw[axis] (0,0) -- (0,1.05) node[midway, right] {\(r\)};
\end{scope}

\draw[axis] (5.0,8.3) -- (6.65,8.3);
\node[lbl] at (5.8,9.0) {\(r\ll R\)};

% ---------------------------------------------------------
% Середина: локальный патч
% ---------------------------------------------------------

\node[lbl] at (0,6.15) {\(\Sigma_{\mathrm{loc}}\)};

\begin{scope}[shift={(-4.0,0.95)}]
    % рамка патча
    \draw[boundary] (0,0) rectangle (8.45,5.35);

    % координаты
    \node[sml] at (4.2,-0.55) {\(z\)};
    \node[sml, rotate=90] at (-0.7,2.65) {\(\vartheta\)};

    % базис
    \draw[axis] (0.95,0.85) -- (2.95,0.85) node[right] {\(e_z\)};
    \draw[axis] (0.95,0.85) -- (0.95,2.85) node[above] {\(e_\vartheta\)};

    % точка и нормаль
    \fill (3.2,2.55) circle (2pt);
    \node[sml, above right] at (3.2,2.55) {\(x\)};
    \draw[axis] (3.2,2.55) -- (4.5,3.95) node[above] {\(n(x)\)};

    % вертикальные стрелки контактов сверху
    \foreach \xx in {1.0,1.9,3.0,4.15,5.45,6.75} {
        \draw[axis] (\xx,6.05) -- (\xx,5.35);
    }
    \node[sml] at (2.75,6.45) {контакты};

    % шнековые линии
    \foreach \c in {-0.8,0.15,1.1,2.05} {
        \draw[screw] (0.45,\c+0.95) -- (4.25,\c+4.05);
    }
    \node[sml] at (5.95,4.25) {\(s_{\rm sh}\)};
    \draw[screw] (5.0,3.8) -- (5.85,4.35);

    % траектории Амеров
    \foreach \yy in {1.25,1.95,2.65,3.35,4.05} {
        \draw[amer] (5.1,\yy) -- (7.15,\yy+0.46);
    }
    \node[sml] at (7.25,2.65) {Амеры};

    % h_A и h_S
    \draw[amer] (3.2,2.55) -- (4.9,2.95);
    \node[sml, below right] at (4.9,2.95) {\(h_A\)};

    \draw[screw] (3.2,2.55) -- (4.65,4.0);
    \node[sml, above right] at (4.65,4.0) {\(h_S\)};

    % краткие поля
    \node[sml, align=left] at (1.95,3.55) {\(\Phi^\pm,\ \lambda\)};
    \node[sml, align=left] at (2.0,2.95) {\(K,\ G,\ A_{ij}\)};
    \node[sml, align=left] at (2.25,1.25) {\(j_{\mathrm{tor}}^{\mathrm{loc}},\ j_{\mathrm{sp}}^{\mathrm{loc}}\)};

    % различение директоров
    \node[sml, align=left] at (6.55,1.0) {\(h_A \neq h_S\)};
    \node[sml, align=left] at (6.65,0.45) {\(e_\parallel \neq h_A,h_S\)};
\end{scope}

% ---------------------------------------------------------
% Низ: минимальный вывод из геометрии
% ---------------------------------------------------------

\node[lbl] at (0,-1.15) {\(\Sigma_{\mathrm{loc}}\subset\Sigma\)};
\node[lbl] at (0,-2.35) {\((e_z,e_\vartheta,n),\ h_A,\ h_S\)};
\node[lbl] at (0,-3.55) {\(r\ll R\ \Rightarrow\) локальная квазицилиндрическая запись};

\end{tikzpicture}%
}
\end{center}

\restoregeometry
\endgroup
\clearpage






% =========================================================
% ВСТАВКА 2: КОМПАКТНАЯ ЛЕГЕНДА / ТАБЛИЦА К РИСУНКУ
% =========================================================

\subsubsection{Легенда к локальному патчу и дискретно-модовой записи}

Для удобства чтения соберём основные обозначения, используемые на рис.~\ref{fig:quasi_cyl_patch_scheme} и в локальных формулах, в компактную таблицу.

\begin{table}[htbp]
\centering
\small
\begin{tabular}{>{\raggedright\arraybackslash}p{2.2cm} >{\raggedright\arraybackslash}p{9.8cm}}
\toprule
Обозначение & Смысл \\
\midrule
\(\Sigma\) & глобальная граничная поверхность ЭВО \\
\(\Sigma_{\mathrm{loc}}\) & локальный квазицилиндрический патч на \(\Sigma\) \\
\(z\) & координата вдоль большой окружности тора \\
\(\vartheta\) & координата по малой окружности тора \\
\(n\) & внешняя единичная нормаль к \(\Sigma_{\mathrm{loc}}\) \\
\(e_z\) & касательное направление вдоль большой окружности (\(e_z\leftrightarrow t_\varphi\)) \\
\(e_\vartheta\) & касательное направление по малой окружности (\(e_\vartheta\leftrightarrow t_\theta\)) \\
\(e_\parallel\) & ось симметрии объекта и ось тестов §17; не совпадает с \(h_A\) и \(h_S\) \\
\(h_A\) & директор траекторий Амеров в канале; задаётся канальной навивкой \(m_A\) \\
\(h_S\) & директор шнековой моды ГП; задаётся модой \(m_S\), хиральностью \(\chi\) и направлением бега \(s_{\rm sh}\) \\
\(\Phi^\pm,\ \lambda\) & интенсивности контактов на ГП \\
\(K,G,A_{ij}\) & прокси-ярлыки локальной статистики контактов \\
\(g_\top,\ a_\top\) & касательные проекции \(G\) и \(A n\) \\
\(m_A\) & квантовое число канальной навивки траекторий Амеров \\
\(m_S\) & базовая шнековая мода ГП \\
\(S1,S2,S4\) & базовые моды шнека: 1, 2 и 4 оборота Амера \\
\(S^\*\) & производная/гармоническая/составная мода \\
\(q_{\mathrm{eff}}\) & только эффективный локальный наклон патча; не главный физический носитель шнека \\
\(j_{\mathrm{tor}}^{\mathrm{loc}}\) & локальный тороидальный вклад внутренней кинетики \\
\(j_{\mathrm{sp}}^{\mathrm{loc}}\) & локальный пространственно-дрейфовый вклад, производный от tor-вклада \\
\(j_{VR}^{\mathrm{int,loc}}\) & полный внутренний вклад на патче \\
\bottomrule
\end{tabular}
\caption{
Компактная легенда к локальному квазицилиндрическому патчу. 
Таблица подчёркивает, что \(\Sigma_{\mathrm{loc}}\) является только локальным языком записи и вычисления, а физически ведущими параметрами шнекового режима служат дискретные величины \(m_A\) и \(m_S\), а не непрерывный параметр наклона.
}
\label{tab:quasi_cyl_patch_legend}
\end{table}

Дополнительно удобно явно выписать локальную последовательность вычисления:
\[
j(x)=\lambda(x)\,\mathbb{E}[\Delta p\mid x],
\]
\[
j_{\mathrm{tor}}^{\mathrm{loc}}
=
\chi\, s_{\mathrm{sh}}\,
\mathcal{T}(m_A;K,g_\top,a_\top,n),
\]
\[
j_{\mathrm{sp}}^{\mathrm{loc}}
=
\beta(K;m_A,m_S)\,
w_{\mathrm{loc}}(x)\,
\Gamma_{A,S}(x)\,
h_S(x),
\qquad
\Gamma_{A,S}(x)=
\bigl(j_{\mathrm{tor}}^{\mathrm{loc}}\cdot h_A\bigr)\,
\Xi(m_A,m_S),
\]
\[
j_{VR}^{\mathrm{int,loc}}
=
j_{\mathrm{tor}}^{\mathrm{loc}}
+
j_{\mathrm{sp}}^{\mathrm{loc}},
\]
после чего локальная запись снова собирается в глобальные интегралы по \(\Sigma\).

\subsection{Почему непрерывный параметр \(q\) оказался недостаточным}

Прежняя непрерывная \(q\)-геометризация была полезным промежуточным языком. Она позволила впервые выразить идею, что пространственный дрейф двутория невозможен без зашитой в геометрию винтовой организации, а также развести тороидальную и пространственную компоненты внутреннего вклада.

Исторически это выглядело так:
\[
h(x)=\frac{t_\varphi(x)+q\,t_\theta(x)}{\|t_\varphi(x)+q\,t_\theta(x)\|},
\qquad
q\in\mathbb{R}\setminus\{0\},
\]
и затем пространственный вклад записывался как производный от tor-вклада через этот директор намотки.

Однако такой язык оказался недостаточным по физическому смыслу. Он позволял параметризовать локальный наклон, но не удерживал квантовую связность шнекового процесса с канальной навивкой. Если шнек действительно обусловлен квантовыми числами каналов и навивкой траекторий Амеров, то одного непрерывного \(q\) уже недостаточно, чтобы не потерять физику процесса.

Поэтому в текущей редакции \(q\) снимается как главный физический носитель шнека. Он допускается лишь как вспомогательный эффективный параметр локального наклона:
\[
q_{\mathrm{eff}},
\]
и не используется как самостоятельное основание физической интерпретации режима.

\subsection{Два касательных директора: \(h_A\) и \(h_S\)}

Ключевое уточнение нового каркаса состоит в том, что на \(\Sigma_{\mathrm{loc}}\) нужно различать два касательных семейства.

\subsubsection{Директор траектории Амеров}

\[
h_A=h_A(m_A;x),
\]
где \(m_A\in\mathbb{Z}_{>0}\) --- дискретное квантовое число канальной навивки. Вектор \(h_A\) задаёт локальное касательное направление, связанное именно с траекторией Амеров в канале.

\subsubsection{Директор шнековой моды ГП}

\[
h_S=h_S(m_S,\chi,s_{\mathrm{sh}};x),
\]
где \(m_S\) задаёт дискретную базовую пространственную моду шнека на ГП. Вектор \(h_S\) описывает не канал как таковой, а шнековую организацию самой граничной поверхности.

\subsubsection{Почему это разведение необходимо}

Шнек не отождествляется с канальной траекторией. Если их не развести, то шнековая геометрия незаметно растворяется в одной непрерывной ``наклонной'' записи и перестаёт быть отличимой как отдельный динамический режим ГП.

Именно поэтому в новом каркасе:
\begin{itemize}
\item \(h_A\) относится к канальной навивке;
\item \(h_S\) относится к шнековой моде ГП;
\item \(e_\parallel\) остаётся осью §17 и не совпадает ни с \(h_A\), ни с \(h_S\).
\end{itemize}

Если на патче используется гипотеза двух винтовых семейств, она должна быть явно помечена как кандидат:
\[
h_A\cdot h_S=0.
\]
При необходимости допускается и кандидат-гипотеза о специальном угле амерной навивки на развёртке, но она не должна быть скрыта внутри параметра.

\subsection{Дискретно-модовая схема шнека: \(m_A\), \(m_S\), \(S1/S2/S4\)}

\subsubsection{Почему шнек теперь читается дискретно}

Теперь физика шнека читается не через произвольный непрерывный наклон, а через согласованную пару
\[
(m_A,m_S),
\]
где:
\[
m_A\in\mathbb{Z}_{>0}
\]
--- квантовое число канальной навивки траекторий Амеров,
а
\[
m_S\in\{1,2,4\}
\]
--- базовая пространственная мода шнека на ГП.

Это выражает новую рабочую норму: шнековая навивка не свободна, а согласована с канальной навивкой.

\subsubsection{Таблица базовых и производных режимов}

\begin{center}
\begin{tabular}{|c|c|c|p{7.2cm}|}
\hline
Обозначение & Параметр & Статус & Смысл \\
\hline
\(S1\) & \(m_S=1\) & базовая мода & шнек растянут на 1 оборот Амера \\
\hline
\(S2\) & \(m_S=2\) & базовая мода & шнек растянут на 2 оборота Амера \\
\hline
\(S4\) & \(m_S=4\) & базовая мода & шнек растянут на 4 оборота Амера \\
\hline
\(S^\*\) & \(m_S\notin\{1,2,4\}\) & только производная & гармоническая/составная мода, допустимая только как производная от \(S1/S2/S4\) \\
\hline
\end{tabular}
\end{center}

Таблица нужна не как голая классификация, а как ответ на физический вопрос: почему именно эти режимы следует считать базовыми. В текущем каркасе это те моды, которые принимаются как первичные пространственные режимы шнековой организации ГП; все остальные допустимы лишь как производные и не должны интерпретироваться как новые базовые формы без отдельного обоснования.

\subsubsection{Локальные вклады в новой записи}

Локальный tor-вклад:
\[
j_{\mathrm{tor}}^{\mathrm{loc}}
=
\chi\, s_{\mathrm{sh}}\,
\mathcal{T}(m_A;K,g_\top,a_\top,n),
\]
где минимальная рабочая запись:
\[
\mathcal{T}(m_A;K,g_\top,a_\top,n)
=
\alpha_1(K,m_A)\,n\times g_\top
+
\alpha_2(K,m_A)\,n\times a_\top.
\]

Локальный sp-вклад:
\[
j_{\mathrm{sp}}^{\mathrm{loc}}
=
\beta(K;m_A,m_S)\,
w_{\mathrm{loc}}(x)\,
\Gamma_{A,S}(x)\,
h_S(x),
\]
где
\[
\Gamma_{A,S}(x)=
\bigl(j_{\mathrm{tor}}^{\mathrm{loc}}\cdot h_A\bigr)\,
\Xi(m_A,m_S).
\]

Полный внутренний вклад:
\[
j_{VR}^{\mathrm{int,loc}}
=
j_{\mathrm{tor}}^{\mathrm{loc}}
+
j_{\mathrm{sp}}^{\mathrm{loc}}.
\]

В этой записи сохраняется главный принцип: без tor-режима дрейфа нет, а пространственный вклад строго производен от тороидального.

\subsection{Возврат от локального патча к глобальному интегралу по \(\Sigma\)}

Локальный патч не завершает картину, а только подготавливает её. Итоговый физический отклик остаётся глобальным:
\[
V=\mu_T\int_{\Sigma} j\,dS,
\qquad
\Omega=\tilde\mu_R\int_{\Sigma}(r-r_0)\times j\,dS,
\qquad
V_\parallel=V\cdot e_\parallel.
\]

Практически это означает, что глобальный интеграл может быть собран как сумма по локальным патчам:
\[
\int_{\Sigma}(\cdots)\,dS
\approx
\sum_{\mathrm{patches}}
\int_{\Sigma_{\mathrm{loc}}}(\cdots)\,dS.
\]

Но именно здесь нужно поставить последнюю витринную оговорку: локальная запись не имеет собственной онтологии. Она не заменяет тор и не даёт права интерпретировать \(\Sigma_{\mathrm{loc}}\) как самостоятельную физическую форму ЭВО. Смысл локального патча --- лишь в том, чтобы яснее показать механизм, который потом снова собирается в глобальный тороидальный отклик.

\subsection{Охранные нормы и интерпретационные границы}

Для витрины полезно завершить раздел коротким блоком ограничений:

\begin{itemize}
\item КУБК относится только к балансу числа контактов на \(\Sigma\).
\item Дрейф и вращение возникают только через анизотропию \(\Delta p\) на \(\Sigma\).
\item Прокси \((K,G,A_{ij})\) --- только ярлыки измерения.
\item \(q_{\mathrm{eff}}\) не является главным физическим носителем шнека.
\item Базовые моды шнека читаются через \(m_S\in\{1,2,4\}\), а не через произвольную непрерывную деформацию.
\item Различение \(h_A\) и \(h_S\) обязательно.
\item Любая локальная запись на \(\Sigma_{\mathrm{loc}}\) должна возвращаться к глобальному интегралу по \(\Sigma\).
\end{itemize}

\subsection{Переходный статус старой \(q\)-схемы}

Для непрерывности изложения полезно сохранить короткое историческое замечание.

\paragraph{Исторический слой.}
Прежняя непрерывная \(q\)-геометризация была полезным промежуточным языком, потому что позволила отделить чисто тороидальный вклад от пространственно-дрейфового и впервые зафиксировала, что дрейф двутория не может рождаться без геометрически зашитой винтовой организации. Однако в дальнейшем оказалось, что одной непрерывной \(q\)-параметризации недостаточно для удержания физики шнекового процесса, так как она не фиксирует жёсткую связь шнека с канальной навивкой и квантовыми числами каналов.

Поэтому в текущей редакции \(q\)-схема остаётся в витрине как исторический переходный слой, а физически ведущей становится дискретно-модовая схема:
\[
(m_A,m_S),\qquad S1/S2/S4,\qquad h_A,\ h_S,\qquad q_{\mathrm{eff}}\ \text{только как вспомогательный параметр.}
\]

% Конец файла
